\section{Dataset and Experimental Setup}
\label{sec:setup}

\paragraph{Tasks.} All data are synthetic and generated on the fly; no external dataset is used.
\begin{itemize}
  \item \textbf{Group word problems (RQ2).} Groups $\mathbb{Z}_2$ (2 tokens), $\mathbb{Z}_3$ (3 tokens) and $S_3$
        (6 tokens). Training sequences have length 64; every position is labelled with the prefix product.
        Training draws $4000\times128 = 512{,}000$ fresh sequences per run.
  \item \textbf{Multi-query associative recall (RQ1).} Following \citet{arora2024zoology}: sequence length 256,
        vocabulary 512 (255 possible keys, 256 possible values), $n\in\{16,32,64\}$ key--value pairs, one query
        per pair. Up to $20{,}000\times64 = 1.28$M fresh sequences per run.
\end{itemize}

\paragraph{Evaluation protocol.} Each run is evaluated on fixed held-out sets of 1024 sequences drawn with a
different random seed from the training stream. For state tracking we evaluate at the training length (64)
and at 256 and 512 (4$\times$ and 8$\times$ longer) to test length generalisation. We report token accuracy
over scored positions, mean $\pm$ standard deviation over three seeds for state tracking. For recall we
report the final accuracy and the number of steps until the evaluation accuracy first exceeded 90\%.

\paragraph{Models.} All models have two layers and $d_{\text{model}}=128$; Table~\ref{tab:models} lists the
configurations. Within each task the recurrent models are matched in total state size (8192 floats) wherever
possible, so differences reflect the transition structure rather than memory size; in the recall study the
state size is varied deliberately.

\begin{table}[htbp]\centering
\resizebox{\linewidth}{!}{\begin{tabular}{llrr}\toprule
Task & Model (mixer configuration) & Parameters & State (floats) \\\midrule
State tracking & Transformer (4 heads, RoPE) & 0.33M & grows with $T$ \\
 & Selective SSM ($N=16$, 8 heads of $P=32$), $a\in(0,1)$ or $(-1,1)$ & 0.41M & 8{,}192 \\
 & Gated DeltaNet (4 heads, $d_k=d_v=32$), $\beta\in(0,1)$ or $(0,2)$ & 0.37M & 8{,}192 \\
 & Rotational SSM ($N=16$, 8 angle pairs) & 0.41M & 8{,}192 \\
 & DeltaProduct, $n_h=2$, $\beta\in(0,1)$ or $(0,2)$ & 0.44M & 8{,}192 \\\midrule
Recall & Attention + short conv (1 head) & 0.20M & grows with $T$ \\
 & Selective SSM, $N=16$ / $N=64$ & 0.28M / 0.30M & 8{,}192 / 32{,}768 \\
 & Gated DeltaNet (2 heads, $d_k=d_v=64$) & 0.23M & 16{,}384 \\\bottomrule
\end{tabular}}
\caption{Model configurations (two layers, $d_{\text{model}}=128$; state-tracking models include a SwiGLU MLP
with ratio 2, recall models have no MLP). Parameter counts include the embedding.}
\label{tab:models}
\end{table}

\paragraph{Hyperparameters.} State tracking: 4000 steps, batch 128, peak learning rate $10^{-3}$, seeds 0--2.
Recall: up to 20{,}000 steps, batch 64, learning rate $10^{-3}$, early stopping at 99\% accuracy, evaluation
every 500 steps, one seed; this budget was chosen by the calibration experiment described in
Section~\ref{sec:discussion}. Common settings are given in Section~\ref{sec:method} (AdamW, warm-up and cosine
schedule, clipping, bfloat16). Efficiency benchmarks use single layers with $d_{\text{model}}=256$, batch 4,
bfloat16; the scan-only benchmark uses float32 with 8 heads and $P=N=64$.

\paragraph{Baselines.} The Transformer is the reference for both abilities: it has an exact memory, so it is
the natural upper bound for recall, and it is known to fail at length-generalising state tracking. Within
each recurrent family the variant with positive eigenvalues is the baseline for the variants that extend the
eigenvalue range or the transition structure.

\paragraph{Hardware and software.} All experiments ran on a single NVIDIA GeForce RTX 4050 Laptop GPU (6\,GB,
Ada architecture) under Ubuntu in WSL2 on Windows, with PyTorch 2 (CUDA build) and Triton. Training one
state-tracking model takes 1.3--5 minutes; the whole study, including the discarded first recall sweep,
took roughly 11 GPU-hours.
